\documentclass[12pt, letterpaper]{article}

\usepackage[utf8]{inputenc}
\usepackage[T1]{fontenc}
\usepackage[spanish]{babel}
\usepackage{geometry}
\usepackage{amsmath, amssymb, amsthm}
\usepackage{booktabs}
\usepackage{array}
\usepackage{microtype}
\usepackage{setspace}
\usepackage{parskip}
\usepackage{titlesec}
\usepackage{fancyhdr}
\usepackage{enumitem}
\usepackage{bm}
\usepackage{mathtools}
\usepackage{tcolorbox}
\usepackage{hyperref}

\tcbuselibrary{skins, breakable}

\geometry{top=2.8cm, bottom=3.0cm, left=3.2cm, right=3.2cm}
\setstretch{1.20}

\pagestyle{fancy}
\fancyhf{}
\fancyhead[C]{\small\textsc{Econometria · Gujarati --- Ejercicio 10.12}}
\fancyfoot[C]{\small\thepage}
\renewcommand{\headrulewidth}{0.4pt}
\renewcommand{\footrulewidth}{0pt}

\titleformat{\section}
  {\large\bfseries\scshape}{\thesection.}{0.6em}{}
  [\vspace{-4pt}\rule{\linewidth}{0.4pt}]
\titleformat{\subsection}{\normalsize\bfseries}{\thesubsection.}{0.5em}{}
\titleformat{\subsubsection}{\normalsize\itshape}{}{0em}{}

\newtheorem{prop}{Proposicion}
\theoremstyle{remark}
\newtheorem{obs}{Observacion}

\newtcolorbox{verdict}[2]{
  colback=white, colframe=black,
  boxrule=0.5pt, arc=3pt,
  left=8pt, right=8pt, top=5pt, bottom=5pt,
  title={\textbf{Inciso #1 --- #2}},
  fonttitle=\small\bfseries,
  breakable
}

\newtcolorbox{clave}{
  colback=white, colframe=black,
  boxrule=0.5pt, arc=3pt,
  left=8pt, right=8pt, top=6pt, bottom=6pt,
  breakable
}

\newcommand{\E}{\mathbb{E}}
\newcommand{\Var}{\operatorname{Var}}
\newcommand{\Cov}{\operatorname{Cov}}
\newcommand{\verdadero}{\textsc{Falso}}
\newcommand{\falso}{\textsc{Falso}}
\newcommand{\incierto}{\textsc{Incierto}}

\begin{document}

\begin{center}
  {\LARGE\bfseries Verdadero, Falso o Incierto:}\\[6pt]
  {\LARGE\bfseries Nueve Afirmaciones sobre Multicolinealidad}\\[12pt]
  {\large\itshape Justificacion formal con algebra y demostraciones}\\[14pt]
  \rule{0.55\linewidth}{0.6pt}\\[8pt]
  {\normalsize Ejercicio~10.12 --- \textit{Econometria}, Gujarati \& Porter, 5.a ed.}\\[4pt]
  {\small\textsc{Con derivaciones matriciales, formulas de varianza y analisis de casos limite}}
\end{center}

\vspace{10pt}

\noindent\textbf{Nota preliminar.}\quad
El modelo de referencia en todos los incisos es:
\begin{equation}
  Y_i = \beta_1 + \beta_2 X_{2i} + \beta_3 X_{3i} + \cdots + \beta_k X_{ki} + u_i
  \label{eq:modelo}
\end{equation}
bajo los supuestos clasicos de MCO: errores $u_i \overset{\text{iid}}{\sim}(0, \sigma^2)$,
regresores no estocasticos y no singularidad de $\mathbf{X'X}$.
La multicolinealidad \emph{perfecta} es el caso limite en que
$\det(\mathbf{X'X}) = 0$; la multicolinealidad \emph{imperfecta} o
\emph{alta} es cuando $\det(\mathbf{X'X}) \approx 0$ pero $> 0$.

\bigskip

%% ── Inciso (a) ────────────────────────────────────────────────────────────
\begin{verdict}{(a)}{FALSO}
  A pesar de la multicolinealidad perfecta, los estimadores de MCO son MELI.
\end{verdict}

\subsection*{Justificacion}

El estimador MCO es $\hat{\boldsymbol\beta} = (\mathbf{X'X})^{-1}\mathbf{X'y}$.
Bajo multicolinealidad perfecta existe $\boldsymbol\lambda \neq \mathbf{0}$ tal que
$\mathbf{X}\boldsymbol\lambda = \mathbf{0}$, lo que implica:
\begin{equation}
  \det(\mathbf{X'X}) = 0
  \label{eq:det_cero}
\end{equation}
La matriz $\mathbf{X'X}$ es singular y \textbf{no tiene inversa}. Por tanto,
el estimador $(\mathbf{X'X})^{-1}\mathbf{X'y}$ \textbf{no existe} como objeto
matematico unico. Hay infinitos vectores $\boldsymbol\beta$ que minimizan la
suma de cuadrados de los residuos.

El teorema de Gauss-Markov establece que el estimador MCO es MELI
(\textit{Best Linear Unbiased Estimator}, BLUE) bajo los supuestos clasicos,
que incluyen explicitamente que $\mathbf{X'X}$ sea invertible (rango completo
de $\mathbf{X}$). Al violarse este supuesto, el teorema no aplica. Los
estimadores individuales $\hat\beta_j$ son indeterminados con varianza infinita,
condicion que es incompatible con ser ``mejor'' en el sentido de minima varianza.

\begin{equation}
  \Var(\hat\beta_j) = \sigma^2 [(\mathbf{X'X})^{-1}]_{jj}
  \xrightarrow{\det(\mathbf{X'X}) \to 0} \infty
  \label{eq:var_inf}
\end{equation}

\textbf{Conclusion}: \textsc{Falso}. Bajo multicolinealidad perfecta, el
estimador MCO no existe, por lo que no puede ser MELI.

\bigskip

%% ── Inciso (b) ────────────────────────────────────────────────────────────
\begin{verdict}{(b)}{INCIERTO}
  En los casos de alta multicolinealidad, no es posible evaluar la significancia
  individual de uno o mas coeficientes de regresion parcial.
\end{verdict}

\subsection*{Justificacion}

El estadistico $t$ del $j$-esimo coeficiente es:
\begin{equation}
  t_j = \frac{\hat\beta_j}{\text{ee}(\hat\beta_j)}
       = \frac{\hat\beta_j}{\hat\sigma/\sqrt{S_{jj}}} \cdot \frac{1}{\sqrt{\text{VIF}_j}}
  \label{eq:t_vif}
\end{equation}
Con VIF$_j$ alto, el denominador se infla y $|t_j|$ se reduce, lo que
\emph{tiende} a producir no rechazo de $H_0: \beta_j = 0$ aunque $\beta_j \neq 0$
(error tipo~II). Sin embargo:

\begin{itemize}[leftmargin=1.6em, itemsep=3pt]
  \item El estadistico $t$ \emph{puede calcularse siempre} que $\mathbf{X'X}$
    sea invertible (colinealidad imperfecta). No es que sea imposible evaluarlo;
    es que el resultado puede ser no significativo por razones de colinealidad,
    no de irrelevancia.
  \item Si el efecto real $\beta_j$ es suficientemente grande o el tamano
    muestral $n$ es suficientemente grande, el coeficiente puede seguir siendo
    significativo incluso con VIF elevado, como se demostro en el ejercicio~10.9.
\end{itemize}

La afirmacion es verdadera en el sentido de que la prueba $t$ \emph{pierde poder}
bajo alta colinealidad, pero es falsa si se interpreta como que la evaluacion
es imposible. La palabra ``no es posible'' hace la afirmacion incorrecta.

\textbf{Conclusion}: \textsc{Incierto}. La prueba $t$ se puede calcular pero
su poder se reduce severamente; puede ser significativa si el efecto real es
grande.

\bigskip

%% ── Inciso (c) ────────────────────────────────────────────────────────────
\begin{verdict}{(c)}{VERDADERO}
  Si una regresion auxiliar muestra que una $R^2_j$ particular es alta, hay
  evidencia clara de alta colinealidad.
\end{verdict}

\subsection*{Justificacion}

La regresion auxiliar de $X_j$ sobre los demas $k-1$ regresores produce un
$R^2_j$. El VIF se define como:
\begin{equation}
  \text{VIF}_j = \frac{1}{1 - R^2_j} \geq 1
  \label{eq:VIF_R2}
\end{equation}

Un $R^2_j$ alto significa que $X_j$ es bien explicada por una combinacion
lineal de los otros regresores, es decir, que $X_j$ es aproximadamente
redundante en el modelo. Concretamente:

\begin{itemize}[leftmargin=1.6em, itemsep=3pt]
  \item $R^2_j = 0.9$ $\Rightarrow$ VIF $= 10$: colinealidad severa.
  \item $R^2_j = 0.99$ $\Rightarrow$ VIF $= 100$: colinealidad muy severa.
  \item $R^2_j \to 1$ $\Rightarrow$ VIF $\to \infty$: aproximacion a multicolinealidad perfecta.
\end{itemize}

La varianza de $\hat\beta_j$ escala con VIF$_j$:
\begin{equation}
  \Var(\hat\beta_j) = \frac{\sigma^2}{S_{jj}} \cdot \text{VIF}_j
  \label{eq:var_VIF}
\end{equation}
Un $R^2_j$ alto implica directamente VIF alto y varianzas infladas. Esta es
la tecnica de diagnostico de Klein (1962): si $R^2_j > R^2$ del modelo principal,
la colinealidad es problematica.

\textbf{Conclusion}: \textsc{Verdadero}. Un $R^2_j$ alto en la regresion
auxiliar es evidencia directa de alta colinealidad de $X_j$ con los demas
regresores.

\bigskip

%% ── Inciso (d) ────────────────────────────────────────────────────────────
\begin{verdict}{(d)}{FALSO}
  Las correlaciones altas entre parejas de regresoras no sugieren una alta
  multicolinealidad.
\end{verdict}

\subsection*{Justificacion}

Las correlaciones de orden cero $r_{jk}$ entre pares de regresores son una
\textbf{condicion suficiente} para la existencia de multicolinealidad.
Si $r_{jk} \approx 1$, la relacion lineal aproximada $X_j \approx aX_k + b$
implica que $\mathbf{X'X}$ se aproxima a la singularidad. En el caso de
dos regresores, el determinante de la submatriz de productos cruzados es:
\begin{equation}
  \det(\mathbf{S}) = S_{22}S_{33} - S_{23}^2
                   = S_{22}S_{33}(1 - r_{23}^2)
  \label{eq:det_S}
\end{equation}
Cuando $r_{23} \to 1$, $\det(\mathbf{S}) \to 0$: multicolinealidad perfecta.
Por tanto, correlaciones altas por pares \emph{si sugieren} alta colinealidad.

La aclaracion importante es que correlaciones bajas \emph{no garantizan} ausencia
de colinealidad: en modelos con $k \geq 3$ regresores puede existir colinealidad
entre combinaciones lineales de varios regresores aunque cada par individual
tenga correlacion baja. Pero eso no cambia la afirmacion: correlaciones altas
\emph{si} sugieren colinealidad, lo cual hace la afirmacion del inciso \emph{falsa}.

\textbf{Conclusion}: \textsc{Falso}. Correlaciones altas entre pares de
regresores son condicion suficiente para la existencia de multicolinealidad.

\bigskip

%% ── Inciso (e) ────────────────────────────────────────────────────────────
\begin{verdict}{(e)}{INCIERTO}
  La multicolinealidad es inofensiva si el objetivo del analisis es solo la
  prediccion.
\end{verdict}

\subsection*{Justificacion}

Para la prediccion, el objeto de interes es $\hat Y_0 = \mathbf{x}_0'\hat{\boldsymbol\beta}$,
no los coeficientes individuales. La varianza del error de prediccion es:
\begin{equation}
  \Var(\hat Y_0 - Y_0) = \sigma^2\left(1 + \mathbf{x}_0'(\mathbf{X'X})^{-1}\mathbf{x}_0\right)
  \label{eq:var_pred}
\end{equation}

La multicolinealidad infla $(\mathbf{X'X})^{-1}$, pero si el vector de
prediccion $\mathbf{x}_0$ cae dentro del espacio de colinealidad de los
datos de entrenamiento, el segundo termino puede seguir siendo pequeno.
Formalmente, si la estructura de colinealidad de los datos futuros es
identica a la de la muestra original, las predicciones son precisas aunque
los coeficientes individuales sean imprecisos.

Sin embargo, la afirmacion es \emph{incierta} porque:
\begin{itemize}[leftmargin=1.6em, itemsep=3pt]
  \item Si $\mathbf{x}_0$ contiene valores de los regresores que rompen la
    estructura de colinealidad del entrenamiento, la prediccion sera muy
    imprecisa.
  \item En la practica, la colinealidad cambia a lo largo del tiempo y los
    nuevos datos rara vez obedecen exactamente la misma estructura lineal.
  \item Incluso para prediccion, los intervalos de prediccion son mas
    amplios bajo alta colinealidad, reduciendo la utilidad practica.
\end{itemize}

\textbf{Conclusion}: \textsc{Incierto}. La multicolinealidad es inofensiva
para la prediccion \emph{solo si} la estructura de colinealidad se mantiene
en los datos de prediccion. En la practica, esto raramente se garantiza.

\bigskip

%% ── Inciso (f) ────────────────────────────────────────────────────────────
\begin{verdict}{(f)}{VERDADERO}
  Entre mayor sea el FIV, ceteris paribus, mas grandes seran las varianzas de
  los estimadores de MCO.
\end{verdict}

\subsection*{Justificacion}

Por definicion del VIF (Factor de Inflacion de la Varianza), la varianza del
estimador MCO de $\beta_j$ es:
\begin{equation}
  \Var(\hat\beta_j)
    = \frac{\sigma^2}{S_{jj}(1-R_j^2)}
    = \frac{\sigma^2}{S_{jj}} \cdot \text{VIF}_j
  \label{eq:var_vif_def}
\end{equation}
donde $S_{jj} = \sum_i x_{ji}^2$ y $R_j^2$ es el $R^2$ de la regresion
auxiliar de $X_j$ sobre los demas regresores. La relacion es \textbf{lineal
y directa}: duplicar el VIF duplica la varianza, todo lo demas constante.
La desigualdad es estricta:
\begin{equation}
  \text{VIF}_j \geq 1 \quad \forall j, \quad
  \text{con igualdad solo si } R_j^2 = 0 \text{ (ortogonalidad)}
  \label{eq:VIF_geq1}
\end{equation}
El VIF cuantifica exactamente cuanto se infla la varianza respecto al
escenario ideal de regresores ortogonales. Un VIF de 10 significa que la
varianza es 10 veces mayor de lo que seria con ortogonalidad.

\textbf{Conclusion}: \textsc{Verdadero}. La varianza de $\hat\beta_j$ es
directamente proporcional a VIF$_j$, por definicion.

\bigskip

%% ── Inciso (g) ────────────────────────────────────────────────────────────
\begin{verdict}{(g)}{FALSO}
  La tolerancia (TOL) es una medida de multicolinealidad mejor que el FIV.
\end{verdict}

\subsection*{Justificacion}

La tolerancia se define como:
\begin{equation}
  \text{TOL}_j = 1 - R_j^2
  \label{eq:TOL}
\end{equation}
y su relacion con el VIF es exactamente:
\begin{equation}
  \text{TOL}_j = \frac{1}{\text{VIF}_j}
  \label{eq:TOL_VIF}
\end{equation}
Dado que TOL y VIF son funciones biunivocas una de la otra
(transformacion monotona decreciente), contienen \textbf{exactamente la
misma informacion}. Conocer uno equivale a conocer el otro:
\begin{align}
  \text{TOL}_j = 0.10 &\iff \text{VIF}_j = 10 \quad \text{(colinealidad severa)}\\
  \text{TOL}_j = 0.01 &\iff \text{VIF}_j = 100 \quad \text{(colinealidad muy severa)}
\end{align}
No existe ninguna situacion en la que TOL revele informacion que VIF no
revele, ni viceversa. Decir que una es ``mejor'' que la otra carece de
sentido matematico. Son dos escales distintas para medir lo mismo.

\textbf{Conclusion}: \textsc{Falso}. TOL y VIF son el inverso multiplicativo
una de la otra. Proveen identica informacion; ninguna es superior.

\bigskip

%% ── Inciso (h) ────────────────────────────────────────────────────────────
\begin{verdict}{(h)}{FALSO}
  No podra obtener un valor $R^2$ elevado en una regresion multiple si todos
  los coeficientes parciales de pendiente no son estadisticamente significativos,
  en lo individual, con base en la prueba $t$ usual.
\end{verdict}

\subsection*{Justificacion}

La afirmacion es \textbf{falsa}: la combinacion de $R^2$ alto con todos los
$t_j$ no significativos es precisamente el \emph{sintoma clasico} de
multicolinealidad severa. La compatibilidad de estas dos situaciones se
demuestra formalmente a continuacion.

El estadistico $F$ del modelo completo contrasta $H_0: \beta_2 = \cdots = \beta_k = 0$:
\begin{equation}
  F = \frac{R^2/(k-1)}{(1-R^2)/(n-k)}
  \label{eq:F_global}
\end{equation}
Un $R^2$ alto implica un $F$ global significativo, lo que rechaza
$H_0$ conjuntamente. Sin embargo, cada estadistico $t$ individual es:
\begin{equation}
  t_j = \frac{\hat\beta_j}{\hat\sigma\sqrt{[(\mathbf{X'X})^{-1}]_{jj}}}
  \label{eq:t_individual}
\end{equation}
Bajo colinealidad, $[(\mathbf{X'X})^{-1}]_{jj}$ es grande (VIF elevado),
inflando el denominador de cada $t_j$ y haciendolo no significativo. La
varianza del error $\hat\sigma^2$ puede ser pequena (buen ajuste global),
pero los errores estandar individuales pueden ser grandes simultaneamente.

Estas dos condiciones son perfectamente compatibles:
\begin{equation}
  R^2 \approx 1 \quad \text{y} \quad |t_j| < t_c \quad \forall j
  \label{eq:paradoja}
\end{equation}
Es la paradoja del $R^2$ alto con $t$'s insignificantes, documentada en
el ejemplo~10.1 de Gujarati con $R^2 = 0.9608$ y ambos $t$'s menores que~2.

\textbf{Conclusion}: \textsc{Falso}. Es completamente posible tener $R^2$
elevado con todos los $t_j$ no significativos. Esa combinacion es de hecho
el indicador mas conspicuo de multicolinealidad severa.

\bigskip

%% ── Inciso (i) ────────────────────────────────────────────────────────────
\begin{verdict}{(i)}{VERDADERO}
  En la regresion de $Y$ sobre $X_2$ y $X_3$, poca variabilidad en $X_3$
  aumenta $\Var(\hat\beta_3)$. En el extremo, si todas las $X_3$ son
  identicas, $\Var(\hat\beta_3)$ seria infinita.
\end{verdict}

\subsection*{Justificacion}

La varianza del estimador MCO de $\beta_3$ en la regresion multiple es:
\begin{equation}
  \Var(\hat\beta_3)
    = \frac{\sigma^2}{\sum x_{3i}^2 \cdot (1 - r_{23}^2)}
    = \frac{\sigma^2}{S_{33}} \cdot \text{VIF}_3
  \label{eq:var_beta3}
\end{equation}
donde $S_{33} = \sum_i (X_{3i} - \bar X_3)^2$ es la varianza muestral
total de $X_3$ (multiplicada por $n-1$). La varianza de $\hat\beta_3$
es \textbf{inversamente proporcional a $S_{33}$}:

\begin{itemize}[leftmargin=1.6em, itemsep=3pt]
  \item Si la variabilidad de $X_3$ es pequena ($S_{33}$ pequeno),
    $\Var(\hat\beta_3)$ es grande: hay poca informacion para estimar $\beta_3$.
  \item En el extremo, si todas las $X_{3i}$ son identicas ($S_{33} = 0$):
    \begin{equation}
      \Var(\hat\beta_3) = \frac{\sigma^2}{0 \cdot (1-r_{23}^2)} \to \infty
      \label{eq:var_inf_x3}
    \end{equation}
    Ademas, si $X_3$ es constante, es perfectamente colineal con la
    constante del modelo (columna de unos), $\det(\mathbf{X'X}) = 0$
    y el estimador no existe.
\end{itemize}

Este fenomeno —variabilidad insuficiente en un regresor— se denomina
\textbf{micronumerosidad} y es conceptualmente analogo a la multicolinealidad.
La intuicion es clara: para estimar el efecto de $X_3$ sobre $Y$ se necesita
variacion en $X_3$. Sin variacion, no hay informacion.

\textbf{Conclusion}: \textsc{Verdadero}. La varianza de $\hat\beta_3$ es
inversamente proporcional a $S_{33}$; con $S_{33} = 0$ la varianza es infinita
y el estimador no existe.

\bigskip

%% ── Resumen ───────────────────────────────────────────────────────────────
\section{Tabla resumen}

\begin{table}[ht]
\centering
\caption{Veredictos del ejercicio~10.12}
\label{tab:resumen}
\small
\begin{tabular}{clp{8.5cm}}
\toprule
\textbf{Inciso} & \textbf{Veredicto} & \textbf{Razon principal} \\
\midrule
(a) & Falso    & Con multicolinealidad perfecta $(\mathbf{X'X})^{-1}$ no existe; el estimador MCO es indeterminado y no puede ser MELI. \\[4pt]
(b) & Incierto & La prueba $t$ se puede calcular pero pierde poder; bajo efectos muy grandes puede ser significativa. \\[4pt]
(c) & Verdadero & $R^2_j$ alto en regresion auxiliar $\Rightarrow$ VIF$_j$ alto $\Rightarrow$ varianza inflada; evidencia directa de colinealidad. \\[4pt]
(d) & Falso    & Correlaciones altas son condicion suficiente (no necesaria) de colinealidad; si la sugieren. \\[4pt]
(e) & Incierto & Inofensiva para prediccion solo si la estructura de colinealidad se mantiene en datos futuros. \\[4pt]
(f) & Verdadero & $\Var(\hat\beta_j) = (\sigma^2/S_{jj})\cdot$VIF$_j$; relacion directa y proporcional por definicion. \\[4pt]
(g) & Falso    & TOL $= 1/$VIF; contienen identica informacion; ninguna es superior. \\[4pt]
(h) & Falso    & $R^2$ alto con $t$'s insignificantes es el sintoma clasico de multicolinealidad. Son perfectamente compatibles. \\[4pt]
(i) & Verdadero & $\Var(\hat\beta_3) \propto 1/S_{33}$; si $S_{33} = 0$ la varianza es infinita y el estimador no existe. \\
\bottomrule
\end{tabular}
\end{table}

\appendix

\section{Demostracion de la compatibilidad de $R^2$ alto con $t$'s insignificantes}

Sea el modelo $Y = \beta_1 + \beta_2 X_2 + \beta_3 X_3 + u$ con
$r_{23} = r$ cercano a~1. Los estadisticos $t$ son:

\begin{equation}
  t_2 = \frac{\hat\beta_2}{\hat\sigma} \cdot \sqrt{\frac{S_{22}(1-r^2)}{1}}
      = \frac{\hat\beta_2}{\hat\sigma/\sqrt{S_{22}}} \cdot \sqrt{1-r^2}
\end{equation}

Cuando $r \to 1$: $\sqrt{1-r^2} \to 0$, por lo que $t_2 \to 0$ aunque
$\hat\beta_2 \neq 0$. El $R^2$ del modelo, en cambio, depende de la variacion
total explicada y puede permanecer alto incluso cuando $r \to 1$, pues las
variables en conjunto siguen explicando bien a $Y$.

\section{Relacion entre varianza de prediccion y colinealidad}

Para el vector de prediccion $\mathbf{x}_0$, la varianza es:
\begin{equation}
  \Var(\hat Y_0) = \sigma^2 \mathbf{x}_0'(\mathbf{X'X})^{-1}\mathbf{x}_0
\end{equation}

Si $\mathbf{x}_0$ es una combinacion lineal de los vectores de entrenamiento
(es decir, cae dentro del espacio columna de $\mathbf{X}$), la varianza puede
ser pequena incluso con colinealidad alta. Pero si $\mathbf{x}_0$ se aleja
de ese espacio, la varianza puede ser enorme. Esto es lo que hace que la
afirmacion del inciso~(e) sea incierta en la practica.

\vspace{14pt}
\noindent\rule{\linewidth}{0.4pt}

\noindent{\small\textit{Nota.} Los resultados de este ejercicio constituyen un
resumen de los conceptos centrales del capitulo~10 de Gujarati \& Porter (2010).
El inciso~(h) reproduce la paradoja del ejemplo~10.1 (consumo, ingreso y
riqueza), donde $R^2 = 0.9608$ y los estadisticos $t$ son $1.14$ y $-0.52$,
ninguno significativo al $5\%$.}

\end{document}
